\chapter{Control Flow}
EDADSL provides several control flow constructs for conditional execution, iteration, and program flow control.

\section{Conditional Execution}
Conditional execution uses the \texttt{if} keyword followed by a boolean expression and a code block. An optional \texttt{else} clause provides an alternative code block executed when the condition is \texttt{False}:
\begin{verbatim}
EBNF:
conditional = "if", boolean_expression, 
              "::", code_block,
              ["else", "::", code_block]

Example:
$flag [bool]= True
if $flag ::{
    f.print("Flag is set")
} else ::{
    f.print("Flag is not set")
}
\end{verbatim}

Basic \texttt{if} without \texttt{else}:
\begin{verbatim}
Example:
$flag [bool]= True
if $flag ::{
    f.print("Flag is set")
}
\end{verbatim}

\section{For Loops}
For loops iterate over a list or set:
\begin{verbatim}
EBNF:
iterator = "for", variable, "E", 
           (list | set), "::", code_block

Example:
$colors [list]= ["red", "green", "blue"]
for $c E $colors ::{
    f.print($c)
}
\end{verbatim}
The loop variable (\texttt{\$c}) takes on each value in the list/set in sequence.

\paragraph{Variable Scoping}
The loop variable is \textbf{not} block-scoped: it persists after the loop exits, retaining the last value assigned during iteration. Loop variables declared at the top level are global. Loop variables declared inside a function are local to that function (see Section~\ref{sec:function-scope}).

\section{Goto Statements}
Goto statements provide unconditional jumps to labeled positions:
\begin{verbatim}
EBNF:
gotomarker = "gotomarker", id
gotojump   = "gotojump", id

Example:
gotojump skip_section
f.print("This will not be printed")
gotomarker skip_section
f.print("Jumped here")
\end{verbatim}
Gotomarkers and gotojumps can cross function and loop boundaries. Nesting is allowed.

\section{End of Program Statements}
\begin{verbatim}
EBNF:
halt = "halt"
noop = "skip"
\end{verbatim}
\begin{itemize}
\item \texttt{halt} -- Terminates program execution and exports all created files
\item \texttt{skip} -- No operation; program continues to next statement
\end{itemize}

\section{Output and Logging}
\begin{verbatim}
EBNF:
echo_statement = "echo", variable
writetolog     = "writetolog", string_expression

Example:
f.print("Hello, EDADSL!")
writetolog "Board design completed successfully"
\end{verbatim}
\begin{itemize}
\item \texttt{echo} -- Keyword that prints a variable reference to the console
\item \texttt{writetolog} -- Keyword that appends a message to the log file
\end{itemize}

\chapter{Functions}
\section{Built-in Functions}

\textbf{Important:} All built-in function calls require the \texttt{f.} prefix. This is defined in the formal grammar (Chapter~3):
\begin{verbatim}
function_call_built_in = "f.", id, "(", 
                        [{ data, "," }, data], ")"
\end{verbatim}
Correct usage: \texttt{f.abs(-3.2)}, \texttt{f.sin(\$x)}, \texttt{f.map(\$list, \$func)}\\
Incorrect: \texttt{abs(-3.2)}, \texttt{sin(\$x)}, \texttt{map(\$list, \$func)}

\texttt{f.print()} takes a \textbf{single} argument. Use \texttt{f.str()} for string concatenation:
\begin{verbatim}
f.print("value = " + $x)              # OK
f.print(f.str("value = ", $x))        # OK
f.print("value = ", $x)               # WRONG - two arguments
\end{verbatim}

\subsection{Mathematical Functions}
\subsubsection{Absolute Value Function}
The Absolute Value Function ( denoted as \texttt{abs} ) returns the Absolute Value of the Argument or the Magnitude of a Vector. \\
Morphology of the Absolute Value Function:
\begin{verbatim}
EBNF:
"f.abs", "(", argument, [ ",", "vec", "=", 
       boolean_data_type ], ")"
\end{verbatim}
If the Type of the Argument does not match the expected Type (a Numeric Value or a List of Numeric Values) a Value Error is raised.\\
Example:
\begin{verbatim}
f.abs(-3.2)
\end{verbatim}
This Example would return a value of \texttt{3.2}.

The Argument may also be a List of Numeric Values. If This is the case the Function is iterated over each Element.\\
Example:
\begin{verbatim}
f.abs([-3, 0.2, 0, -7])
\end{verbatim}
This Example would return \texttt{[3, 0.2, 0, 7]}. \vspace{5 mm}\\
If \texttt{vec = True} is passed a List shall be passed as the Argument and the Function will return Magnitude of the List as an n-dimensional Vector, where n is the Length of the List.\\ %$\texttt{abs}(\vec{v}\texttt{, vec = True}) = \sqrt{\sum_{i=1}^nv_i^2}$. 
Example:
\begin{verbatim}
f.abs( [ 3, -4 ], vec = True )
\end{verbatim}
This Example would return \texttt{5.0}.\\
There are no Domain-Restrictions for the function inside $\mathbb{R}^n$.

\paragraph{Complex Overload}
When the Argument is a Complex Number $z = a + bi$, \texttt{abs} returns the Modulus (Magnitude) of the Complex Number, defined as:
\[
\lvert a + bi \rvert = \sqrt{a^2 + b^2}
\]
The Result is always a Real Number, even if the Argument is Complex.\\
Domain: All Complex Numbers ($\mathbb{C}$).\\
Example:
\begin{verbatim}
f.abs(3 + 4i)   #c# Returns 5.0
f.abs(1 + 1i)   #c# Returns approximately 1.41421
f.abs(0 - 5i)   #c# Returns 5.0
\end{verbatim}

\subsubsection{Parallel Equivalent Function}
The Parallel Equivalent Function ( denoted as \texttt{parallel} ) computes the parallel combination of a list of values. This is the standard formula for parallel resistances, conductances, or any set of admittances. \\
Morphology of the Parallel Equivalent Function:
\begin{verbatim}
EBNF:
"f.parallel", "(", argument, ")"
\end{verbatim}
If the Type of the Argument does not match the expected Type (a List of Numeric Values) a Value Error is raised.\\
The Argument must be a List of Numeric Values (\texttt{int}, \texttt{real}, or \texttt{complex}). The function computes:
\[
\text{parallel}([v_1, v_2, \ldots, v_n]) = \frac{1}{\frac{1}{v_1} + \frac{1}{v_2} + \cdots + \frac{1}{v_n}}
\]
If any Element in the List is Zero, the Function immediately Returns Zero.\\
If any Element is Infinite, it is Ignored (does not contribute to the Result).\\
The Return Type is determined by the Numerical Value of the Result:
\begin{itemize}
\item If the Result is Numerically Equal to an Integer, Return \texttt{int}
\item Otherwise, if the Result is Numerically Equal to a Real, Return \texttt{real}
\item Otherwise, Return \texttt{complex}
\end{itemize}
Examples:
\begin{verbatim}
f.parallel([100, 100])         #c# Returns 50 (int)
f.parallel([100, 200, 300])    #c# Returns 54 (int, softint)
f.parallel([1e3, 2.2e3])       #c# Returns 687.5 (real)
f.parallel([100, 0])           #c# Returns 0 (any zero)
f.parallel([100, 200, 0, 300]) #c# Returns 0 (any zero)
f.parallel([100, inf])           #c# Returns 100 (inf ignored)
f.parallel([100, 200, inf])      #c# Returns 66.666... (inf ignored)
\end{verbatim}

\paragraph{Complex Overload}
When the List contains Complex Numbers, the function computes the parallel combination using Complex Arithmetic. If the Result has a Zero Imaginary Part, it is Returned as \texttt{real}. \\
Example:
\begin{verbatim}
f.parallel([100 + 0i, 200 + 0i])
    #c# Returns 66.666... (real, imag=0)
f.parallel([3 + 4i, 1 - 2i])           #c# Returns complex result
\end{verbatim}
Domain: All Lists of Numeric Values. Empty Lists raise a Value Error. Lists with a single Element return that Element unchanged (after \texttt{softint}(\texttt{softreal}())).\\

\subsubsection{Signum Function}
The Signum Function ( denoted as \texttt{sign} ) returns the Sign of the Argument. \\
Morphology of the Signum Function:
\begin{verbatim}
EBNF:
"f.sign", "(", argument, ")"
\end{verbatim}
If the Type of the Argument does not match the expected Type (a Numeric Value or a List of Numeric Values) a Value Error is raised.\\
Example:
\begin{verbatim}
f.sign(-3.2)
\end{verbatim}
This Example would return a value of \texttt{-1}.

The Argument may also be a List of Numeric Values. If This is the case the Function is iterated over each Element.\\
Example:
\begin{verbatim}
f.sign([-3, 0.2, 0, -7])
\end{verbatim}
This Example would return \texttt{[-1, 1, 0, -1]}.\\
There are no Domain-Restrictions for the function inside $\mathbb{R}^n$.


\subsubsection{Regular Trigonometric Functions}
Important: Angles are interpreted as Radians by Default.\\
The Regular Trigonometric Functions (Sine, Cosine and Tangent) ( denoted as \texttt{sin, cos, tan} ) return the Sine, Cosine, or Tangent of the Argument. \\
Morphology of the Trigonometric Functions:
\begin{verbatim}
EBNF:
( "f.sin" | "f.cos" | "f.tan" ), "(", argument, 
  [ ",", "unit", "=", unit ], ")"

unit = ( "'", ("rad" | "deg" ), "'" )
     | ( '"', ("rad" | "deg" ), '"' )
\end{verbatim}
If the Type of the Argument does not match the expected Type (a Numeric Value or a List of Numeric Values) a Value Error is raised.\\
Example:
\begin{verbatim}
f.sin(1.570796327)
\end{verbatim}
This Example would return a value of around \texttt{1}.

The Argument may also be a List of Numeric Values. If This is the case the Function is iterated over each Element.\\
Example:
\begin{verbatim}
f.sin([0, 30, 90], unit = "deg" )
\end{verbatim}
This Example would return \texttt{[0, 0.5, 1]}.\\
For Sine and Cosine there are no Domain-Restrictions for the functions inside $\mathbb{R}^n$. \\
For Tangent there are Domain-Restrictions inside $\mathbb{R}^n$.\\
It is only defined for $\{x \in \mathbb{R}: \cos x \neq 0 \}^n$

\paragraph{Complex Overload}
All three functions (\texttt{sin, cos, tan}) are also defined for Complex Arguments via Analytic Continuation. For complex $z = a + bi$, the functions are computed using Euler's Formula:
\[
\sin(a + bi) = \sin(a)\cosh(b) + i\cos(a)\sinh(b)
\]
\[
\cos(a + bi) = \cos(a)\cosh(b) - i\sin(a)\sinh(b)
\]
\[
\tan(a + bi) = \frac{\sin(a + bi)}{\cos(a + bi)}
\]
Domain: All Complex Numbers ($\mathbb{C}$) for \texttt{sin} and \texttt{cos}. For \texttt{tan}, all Complex Numbers where $\cos z \neq 0$.\\
Examples:
\begin{verbatim}
f.sin(1 + 2i)        #c# Returns approximately 3.1658 - 1.9596i
f.cos(1 + 2i)        #c# Returns approximately 2.0327 + 3.0519i
f.tan(0 + 1i)        #c# Returns approximately 0.0 + 0.7616i
\end{verbatim}

\subsubsection{Regular Inverse Trigonometric Functions}
Important: Angles are interpreted as Radians by Default.\\
The Regular Inverse Trigonometric Functions (Arcsine, Arccosine and Arctangent) ( denoted as \texttt{asin, acos, atan}, \texttt{arsin, arcos, artan} or \texttt{arcsin, arccos, arctan} ) return the Inverse for the Sine, Cosine, or Tangent of the Argument. \\
Morphology of the Trigonometric Functions:
\begin{verbatim}
EBNF:
inverses_trig_functions, "(", argument, 
  [ ",", "unit", "=", unit ], ")"

inverses_trig_functions = 
    "f.asin" | "f.acos" | "f.atan" 
  | "f.arsin" | "f.arcos" | "f.artan"
  | "f.arcsin" | "f.arccos" | "f.arctan"

unit = ( "'", ("rad" | "deg" ), "'" )
     | ( '"', ("rad" | "deg" ), '"' )
\end{verbatim}
If the Type of the Argument does not match the expected Type (a Numeric Value or a List of Numeric Values) a Value Error is raised.\\
Example:
\begin{verbatim}
f.asin(1)
\end{verbatim}
This Example would return a value of around \texttt{1.570796327}.

The Argument may also be a List of Numeric Values. If This is the case the Function is iterated over each Element.\\
Example:
\begin{verbatim}
f.arcsin([0, 0.5, 1], unit = "deg" )
\end{verbatim}
This Example would return \texttt{[0, 30, 90]}.\\
For Arctangent there is no Domain-Restrictions for the functions inside $\mathbb{R}^n$. \\
For Arcsine and Arccosine the Domain is Restricted to $[-1,1]^n$.\\

\subsubsection{Inverse Tangent Function based on a Vector Argument}

\newpage
\subsubsection{Hyperbolic Functions}
The Hyperbolic Functions (Hyperbolic Sine, Hyperbolic Cosine and Hyperbolic Tangent) ( denoted as \texttt{sinh, cosh, tanh} ) return the Hyperbolic Sine, Hyperbolic Cosine, or Hyperbolic Tangent of the Argument. \\
Morphology of the Hyperbolic Functions:
\begin{verbatim}
EBNF:
( "f.sinh" | "f.cosh" | "f.tanh" ), "(", argument, ")"
\end{verbatim}
If the Type of the Argument does not match the expected Type (a Numeric Value or a List of Numeric Values) a Value Error is raised.\\
Example:
\begin{verbatim}
f.sinh(1.443635475)
\end{verbatim}
This Example would return a value of around \texttt{2}.

The Argument may also be a List of Numeric Values. If This is the case the Function is iterated over each Element.\\
Example:
\begin{verbatim}
f.sinh([0, 1, 2] )
\end{verbatim}
This Example would return around \texttt{[0, 1.175, 3.627]}.\\
There are no Domain-Restrictions for the functions inside $\mathbb{R}^n$.

\paragraph{Complex Overload}
All three functions (\texttt{sinh, cosh, tanh}) are also defined for Complex Arguments. For complex $z = a + bi$, the functions are computed as:
\[
\sinh(a + bi) = \sinh(a)\cos(b) + i\cosh(a)\sin(b)
\]
\[
\cosh(a + bi) = \cosh(a)\cos(b) + i\sinh(a)\sin(b)
\]
\[
\tanh(a + bi) = \frac{\sinh(a + bi)}{\cosh(a + bi)}
\]
Equivalently, the identities $\sinh(z) = -i \sin(iz)$, $\cosh(z) = \cos(iz)$, and $\tanh(z) = \sinh(z) / \cosh(z)$ hold for Complex Arguments.\\
Domain: All Complex Numbers ($\mathbb{C}$) for \texttt{sinh} and \texttt{cosh}. For \texttt{tanh}, all Complex Numbers where $\cosh z \neq 0$.\\
Examples:
\begin{verbatim}
f.sinh(1 + 2i)   #c# Returns approximately -0.4891 + 0.9889i
f.cosh(1 + 2i)   #c# Returns approximately -0.4891 + 0.9889i
f.tanh(0 + 1i)   #c# Returns approximately 0.0 + 0.5524i
\end{verbatim}

\subsubsection{Hyperbolic Inverse Functions}
The Hyperbolic Inverse Functions (Hyperbolic Arcsine, Hyperbolic Arccosine and Hyperbolic Arctangent) ( denoted as \texttt{asinh, acosh, atanh}, \texttt{arsinh, arcosh, artanh} or \texttt{arcsinh, arccosh, arctanh} ) return the Inverse for the Hyperbolic Sine, Hyperbolic Cosine, or Hyperbolic Tangent of the Argument. \\
Morphology of the Hyperbolic Functions:
\begin{verbatim}
EBNF:
inverses_hypb_functions, "(", argument, ")"

inverses_hypb_functions = 
    "f.asinh" | "f.acosh" | "f.atanh" 
  | "f.arsinh" | "f.arcosh" | "f.artanh"
  | "f.arcsinh" | "f.arccosh" | "f.arctanh"
\end{verbatim}
If the Type of the Argument does not match the expected Type (a Numeric Value or a List of Numeric Values) a Value Error is raised.\\
Example:
\begin{verbatim}
f.asinh(1)
\end{verbatim}
This Example would return a value of around \texttt{0.881373587}.

The Argument may also be a List of Numeric Values. If This is the case the Function is iterated over each Element.\\
Example:
\begin{verbatim}
f.arcsinh( [0, 0.5, 1] )
\end{verbatim}
This Example would return \texttt{[0, 0.481, 0.881]}.\\
For the Hyperbolic Arcsine there is no Domain-Restrictions for the function inside $\mathbb{R}^n$. \\
For Hyperbolic Arctangent and Hyperbolic Arccosine the Domain is Restricted to $(-1,1)^n$ and $[1,\infty)^n$ respectively.\\

\paragraph{Complex Overload}
All three functions (\texttt{asinh, acosh, atanh}) are also defined for Complex Arguments using the Complex Logarithm. On the Complex Plane, the Real-Domain Restrictions no longer apply:
\[
\operatorname{asinh}(z) = \ln\!\left(z + \sqrt{z^2 + 1}\right)
\]
\[
\operatorname{acosh}(z) = \ln\!\left(z + \sqrt{z^2 - 1}\right)
\]
\[
\operatorname{atanh}(z) = \frac{1}{2} \ln\!\left(\frac{1 + z}{1 - z}\right)
\]
Domain: All Complex Numbers ($\mathbb{C}$).\\
Examples:
\begin{verbatim}
f.asinh(1 + 2i)   #c# Returns approximately 1.4694 + 1.0987i
f.acosh(1 + 2i)   #c# Returns approximately 1.5286 + 1.1437i
f.atanh(0.5 + 0i) #c# Returns approximately 0.5493 + 0.0i
\end{verbatim}

\subsubsection{Exponential Functions and Logarithms}
\paragraph{Exponential Function}
The \texttt{exp} function returns $e$ raised to the power of the argument. \\
Morphology:
\begin{verbatim}
EBNF:
"f.exp", "(", argument, ")"
\end{verbatim}
Example:
\begin{verbatim}
f.exp(0)    #c# Returns 1.0
f.exp(1)    #c# Returns approximately 2.71828
\end{verbatim}
\paragraph{Complex Overload}
When the Argument is a Complex Number $z = a + bi$, the function is computed using Euler's Formula:
\[
e^{a+bi} = e^a \left(\cos b + i \sin b\right)
\]
In particular, $e^{i\pi} = -1$ (Euler's Identity) and $|e^{a+bi}| = e^a$.\\
Domain: All Complex Numbers ($\mathbb{C}$).\\
Examples:
\begin{verbatim}
f.exp(0)              #c# Returns 1.0
f.exp(1 + 3.14159i)   #c# Returns approximately -2.7183 + 0.0i
f.exp(0 + 3.14159i)
    #c# Returns approximately -1 + 0.0i (Euler's Identity)
f.exp(1 + 1i)         #c# Returns approximately 1.4687 + 2.2874i
\end{verbatim}

\paragraph{Square Root Function}
The \texttt{sqrt} function returns the square root of the argument. \\
Morphology:
\begin{verbatim}
EBNF:
"f.sqrt", "(", argument, ")"
\end{verbatim}
Example:
\begin{verbatim}
f.sqrt(9)   #c# Returns 3.0
f.sqrt(2)   #c# Returns approximately 1.41421
\end{verbatim}
\paragraph{Complex Overload}
Domain for Real Arguments: $[0, \infty)$. When the Argument is a Negative Real or Complex Number, the Principal Square Root is returned, defined such that:
\[
-\pi < \arg(\sqrt{z}) \leq \pi
\]
For $z = r e^{i\theta}$, the Principal Square Root is $\sqrt{z} = \sqrt{r} \, e^{i\theta/2}$.\\
Domain: All Complex Numbers ($\mathbb{C}$), using the Principal Branch.\\
Examples:
\begin{verbatim}
f.sqrt(-1)         #c# Returns 0 + 1i
f.sqrt(-4)         #c# Returns 0 + 2i
f.sqrt(4i)         #c# Returns approximately 1.41421 + 1.41421i
f.sqrt(3 + 4i)    #c# Returns approximately 2 + 1i
f.sqrt(0)          #c# Returns 0.0
\end{verbatim}

\paragraph{Power Function}
The \texttt{pow} function returns the first argument raised to the power of the second argument. \\
Morphology:
\begin{verbatim}
EBNF:
"f.pow", "(", argument, ",", argument, ")"
\end{verbatim}
Example:
\begin{verbatim}
f.pow(2, 3)   #c# Returns 8.0
f.pow(9, 0.5) #c# Returns 3.0
\end{verbatim}
\paragraph{Complex Overload}
The function also accepts Complex Arguments. For complex $z_1$ and $z_2$, the result is computed via:
\[
z_1^{z_2} = e^{z_2 \ln z_1}
\]
where $\ln$ is the Principal Branch of the Complex Logarithm.\\
Domain: All Complex Numbers ($\mathbb{C} \times \mathbb{C}$), using the Principal Branch of $\ln$. Note that $0^z$ is only defined for $\operatorname{Re}(z) > 0$.\\
Examples:
\begin{verbatim}
f.pow(2, 1 + 1i)    #c# Returns approximately 1.538 + 0.999i
f.pow(1 + 1i, 2)    #c# Returns 0 + 2i
f.pow(0 + 1i, 0+1i)
    #c# Returns approximately 0.2079 + 0.0i (i^i = e^{-pi/2})
\end{verbatim}

\paragraph{Natural Logarithm}
The \texttt{ln} function returns the natural logarithm (base $e$) of the argument. \\
Morphology:
\begin{verbatim}
EBNF:
"f.ln", "(", argument, ")"
\end{verbatim}
Example:
\begin{verbatim}
f.ln(1)      #c# Returns 0.0
f.ln(e)      #c# Returns 1.0
\end{verbatim}
\paragraph{Complex Overload}
Domain for Real Arguments: $(0, \infty)$. When the Argument is a Complex Number, the Complex Logarithm (Principal Branch) is returned:
\[
\ln(a + bi) = \ln\lvert a + bi \rvert + i \arg(a + bi)
\]
where $-\pi < \arg(a + bi) \leq \pi$.\\
Domain: All Complex Numbers except $0$ ($\mathbb{C} \setminus \{0\}$).\\
Examples:
\begin{verbatim}
f.ln(1)          #c# Returns 0.0 + 0.0i
f.ln(e)          #c# Returns 1.0 + 0.0i
f.ln(-1)         #c# Returns approximately 0 + 3.1416i
f.ln(1 + 1i)    #c# Returns approximately 0.3466 + 0.7854i
f.ln(0 + 1i)    #c# Returns approximately 0 + 1.5708i (i*pi/2)
\end{verbatim}

\paragraph{Base-10 Logarithm}
The \texttt{log} function returns the base-10 logarithm of the argument. \\
Morphology:
\begin{verbatim}
EBNF:
"f.log", "(", argument, ")"
\end{verbatim}
Example:
\begin{verbatim}
f.log(100)   #c# Returns 2.0
f.log(1)     #c# Returns 0.0
\end{verbatim}
\paragraph{Complex Overload}
Domain for Real Arguments: $(0, \infty)$. When the Argument is a Complex Number, the Complex Logarithm to base 10 is returned:
\[
\log(z) = \frac{\ln(z)}{\ln(10)}
\]
where $\ln$ is the Principal Branch of the Complex Logarithm.\\
Domain: All Complex Numbers except $0$ ($\mathbb{C} \setminus \{0\}$).\\
Examples:
\begin{verbatim}
f.log(100)        #c# Returns 2.0 + 0.0i
f.log(1)          #c# Returns 0.0 + 0.0i
f.log(-100)       #c# Returns approximately 2 + 1.3644i
f.log(10i)        #c# Returns approximately 1 + 1.3644i
\end{verbatim}

\label{sec:complex_functions}
\subsubsection{Complex Functions}
Functions for creating and decomposing complex numbers.

\paragraph{Real Part}
The \texttt{real} function returns the real part of a complex number.
\begin{verbatim}
EBNF:
"f.real", "(", complex, ")"
\end{verbatim}
\begin{verbatim}
f.real(3 + 4i)   #c# Returns 3.0
f.real(-2i)      #c# Returns 0.0
f.real(5)        #c# Returns 5.0 (int promoted)
\end{verbatim}

\paragraph{Imaginary Part}
The \texttt{imag} function returns the imaginary part of a complex number (as a real).
\begin{verbatim}
EBNF:
"f.imag", "(", complex, ")"
\end{verbatim}
\begin{verbatim}
f.imag(3 + 4i)   #c# Returns 4.0
f.imag(-2i)      #c# Returns -2.0
f.imag(5)        #c# Returns 0.0 (real has no imaginary part)
\end{verbatim}

\paragraph{Angle}
The \texttt{angle} function returns the angle (argument) of a complex number in radians.
\begin{verbatim}
EBNF:
"f.angle", "(", complex, ")"
\end{verbatim}
\begin{verbatim}
f.angle(1 + 0i)      #c# Returns 0.0
f.angle(0 + 1i)      #c# Returns 1.5708 (pi/2)
f.angle(-1 + 0i)     #c# Returns 3.14159 (pi)
f.angle(1 + 1i)      #c# Returns 0.7854 (pi/4)
\end{verbatim}

\paragraph{Complex Conjugate}
The \texttt{conj} function returns the complex conjugate: \texttt{conj(a + bi) = a - bi}.
\begin{verbatim}
EBNF:
"f.conj", "(", complex, ")"
\end{verbatim}
\begin{verbatim}
f.conj(3 + 4i)   #c# Returns 3 - 4i
f.conj(-2i)      #c# Returns 2i
f.conj(5)        #c# Returns 5 (real numbers are self-conjugate)
\end{verbatim}

\paragraph{Polar Conversion}
The \texttt{polar} function converts a complex number to polar form, returning a list \texttt{[magnitude, angle]}.
\begin{verbatim}
EBNF:
"f.polar", "(", complex, ")"
\end{verbatim}
\begin{verbatim}
f.polar(3 + 4i)     #c# Returns [5.0, 0.9273]
f.polar(1 + 0i)     #c# Returns [1.0, 0.0]
f.polar(0 + 1i)     #c# Returns [1.0, 1.5708]
\end{verbatim}

\paragraph{Rectangular Conversion}
The \texttt{rect} function creates a complex number from polar form (magnitude and angle in radians).
\begin{verbatim}
EBNF:
"f.rect", "(", real, ",", real, ")"
\end{verbatim}
\begin{verbatim}
f.rect(5, 0.9273)      #c# Returns approximately 3 + 4i
f.rect(1, 0)           #c# Returns 1 + 0i
f.rect(1, 1.5708)      #c# Returns approximately 0 + 1i
\end{verbatim}

\paragraph{Radians to Degrees}
The \texttt{rad2deg} function converts an angle from radians to degrees.
\begin{verbatim}
EBNF:
"f.rad2deg", "(", real, ")"
\end{verbatim}
\begin{verbatim}
f.rad2deg(3.14159)   #c# Returns 180.0
f.rad2deg(1.5708)    #c# Returns 90.0
f.rad2deg(0)         #c# Returns 0.0
\end{verbatim}

\paragraph{Degrees to Radians}
The \texttt{deg2rad} function converts an angle from degrees to radians.
\begin{verbatim}
EBNF:
"f.deg2rad", "(", real, ")"
\end{verbatim}
\begin{verbatim}
f.deg2rad(180)   #c# Returns 3.14159
f.deg2rad(90)    #c# Returns 1.5708
f.deg2rad(0)     #c# Returns 0.0
\end{verbatim}

\subsubsection{Rounding Functions}
\paragraph{Floor Function}
The \texttt{floor} function rounds the argument down to the nearest integer. \\
Morphology:
\begin{verbatim}
EBNF:
"f.floor", "(", argument, ")"
\end{verbatim}
Example:
\begin{verbatim}
f.floor(2.7)   #c# Returns 2
f.floor(-1.3)  #c# Returns -2
\end{verbatim}

\paragraph{Ceiling Function}
The \texttt{ceil} function rounds the argument up to the nearest integer. \\
Morphology:
\begin{verbatim}
EBNF:
"f.ceil", "(", argument, ")"
\end{verbatim}
Example:
\begin{verbatim}
f.ceil(2.3)    #c# Returns 3
f.ceil(-2.7)   #c# Returns -2
\end{verbatim}

\paragraph{Round Function}
The \texttt{round} function rounds the argument to the nearest integer. \\
Morphology:
\begin{verbatim}
EBNF:
"f.round", "(", argument, ")"
\end{verbatim}
Example:
\begin{verbatim}
f.round(2.5)   #c# Returns 3
f.round(2.4)   #c# Returns 2
\end{verbatim}

\subsubsection{Power and Modular Exponentiation Function}
The \texttt{pow} function computes exponentiation and optional modular exponentiation. \\
Morphology:
\begin{verbatim}
EBNF:
"f.pow", "(", argument, ",", argument, 
  [ ",", argument ], ")"
\end{verbatim}
With two arguments, \texttt{pow(a, b)} returns $a^b$. With three arguments, \texttt{pow(a, b, c)} returns $a^b \bmod c$ (modular exponentiation). Modular exponentiation is efficient for large exponents and is used in cryptographic and hardware applications. \\
Examples:
\begin{verbatim}
f.pow(2, 10)        #c# Returns 1024
f.pow(3, 4, 7)      #c# Returns 4  (81 mod 7)
f.pow(2, 32, 100)   #c# Returns 96  (4294967296 mod 100)
f.pow(3, -1, 7)     #c# Returns 5  (modular inverse: 3*5=15==1 mod 7)
f.pow(3, -2, 7)     #c# Returns 4  (5^2=25==4 mod 7, i.e. 3^(-2) mod 7)
f.pow(3, -3, 7)     #c# Returns 6  (5^3=125==6 mod 7, i.e. 3^(-3) mod 7)
f.pow(5, -1, 13)    #c# Returns 8  (5*8=40==1 mod 13)
\end{verbatim}
Domain restrictions: When \texttt{c} is provided, \texttt{a} and \texttt{c} must be \texttt{int} and \texttt{c} must be positive. Without \texttt{c}, \texttt{b} must be non-negative (negative exponents on \texttt{int} would produce fractions). With \texttt{c}, negative \texttt{b} computes $a^b \bmod c$ via the modular multiplicative inverse: first find $a^{-1} \bmod c$ (which exists when $\gcd(a, c) = 1$), then raise to $|b|$.

\subsubsection{Extrema Functions}
\paragraph{Minimum Function}
The \texttt{min} function returns the minimum value in a list. For ndarrays, an optional \texttt{axis} parameter reduces along that axis. \\
Morphology:
\begin{verbatim}
EBNF:
"f.min", "(", list, [ ",", "axis", "=", int ], ")"
\end{verbatim}
Examples:
\begin{verbatim}
f.min([3, 1, 4, 1, 5])              #c# Returns 1
f.min([[3,1],[4,2]], axis=0)        #c# Returns [3, 1]
f.min([[3,1],[4,2]], axis=1)        #c# Returns [1, 2]
\end{verbatim}

\paragraph{Maximum Function}
The \texttt{max} function returns the maximum value in a list. For ndarrays, an optional \texttt{axis} parameter reduces along that axis. \\
Morphology:
\begin{verbatim}
EBNF:
"f.max", "(", list, [ ",", "axis", "=", int ], ")"
\end{verbatim}
Examples:
\begin{verbatim}
f.max([3, 1, 4, 1, 5])              #c# Returns 5
f.max([[3,1],[4,2]], axis=0)        #c# Returns [4, 2]
f.max([[3,1],[4,2]], axis=1)        #c# Returns [3, 4]
\end{verbatim}

\subsubsection{Inverse Tangent Function based on a Vector Argument}
The \texttt{atan2} function computes the arctangent of $y/x$ using the signs of both arguments to determine the correct quadrant. \\
Morphology:
\begin{verbatim}
EBNF:
"f.atan2", "(", argument, ",", argument, ")"
\end{verbatim}
Example:
\begin{verbatim}
f.atan2(1, 0)    #c# Returns approximately 1.5708 (pi/2)
f.atan2(0, 1)    #c# Returns 0.0
\end{verbatim}

\subsubsection{Norm and Normalization Functions}
EDADSL provides functions for computing vector norms and normalizing vectors.

\paragraph{Norm Function}
The \texttt{norm} function computes the p-norm of a vector (list of ints or reals). \\
Morphology of the Norm Function:
\begin{verbatim}
EBNF:
"f.norm", "(", argument, [ ",", "p", "=", 
       real_data_type ], ")"
\end{verbatim}
Parameters:
\begin{itemize}
\item \texttt{argument} -- A list of int or real values (the vector)
\item \texttt{p} -- (Optional) The order of the norm. Defaults to 2 (Euclidean norm).
\end{itemize}

Supported norms:
\begin{itemize}
\item \texttt{p = 1} -- L1 norm (Manhattan distance): $\|\mathbf{x}\|_1 = \sum_{i=1}^n |x_i|$
\item \texttt{p = 2} -- L2 norm (Euclidean distance): $\|\mathbf{x}\|_2 = \sqrt{\sum_{i=1}^n x_i^2}$
\item \texttt{p > 2} -- Lp norm: $\|\mathbf{x}\|_p = \left(\sum_{i=1}^n |x_i|^p\right)^{1/p}$
\item \texttt{p = 0} -- L0 "norm" (count of non-zero elements): $\|\mathbf{x}\|_0 = |\{i : x_i \neq 0\}|$
\item \texttt{p = -1} -- L-infinity norm (Chebyshev distance): $\|\mathbf{x}\|_\infty = \max_i |x_i|$
\end{itemize}

Examples:
\begin{verbatim}
f.norm([3, 4])            #c# Returns 5.0 (L2 norm)
f.norm([3, 4], p=1)       #c# Returns 7 (L1 norm)
f.norm([3, 4], p=0)       #c# Returns 2 (L0 norm)
f.norm([-3, 4, -2], p=-1) #c# Returns 4 (L-inf norm)
\end{verbatim}

If the argument is not a list of ints or reals, a Value Error is raised.

\paragraph{Normalize Function}
The \texttt{normalize} function returns a unit vector in the direction of the input vector. \\
Morphology of the Normalize Function:
\begin{verbatim}
EBNF:
"f.normalize", "(", argument, [ ",", "p", "=", 
              real_data_type ], ")"
\end{verbatim}
Parameters:
\begin{itemize}
\item \texttt{argument} -- A list of int or real values (the vector to normalize)
\item \texttt{p} -- (Optional) The order of the norm to use for normalization. Defaults to 2.
\end{itemize}

The function computes: $\hat{\mathbf{x}} = \frac{\mathbf{x}}{\|\mathbf{x}\|_p}$

Examples:
\begin{verbatim}
f.normalize([3, 4])        #c# Returns [0.6, 0.8] (L2)
f.normalize([3, 4], p=1)   #c# Returns [0.43, 0.57] (L1)
f.normalize([6, 8], p=2)   #c# Returns [0.6, 0.8] (L2)
\end{verbatim}

If the norm of the vector is zero, a Division by Zero Error is raised.

\paragraph{Relationship between abs(vec=True) and norm}
The function \texttt{abs(v, vec=True)} is equivalent to \texttt{norm(v)} (L2 norm by default).


\subsection{String Functions}

\subsubsection{String Conversion}
The \texttt{str} function converts its arguments to a string. \\
Morphology:
\begin{verbatim}
EBNF:
"f.str", "(", { argument }, ")"
\end{verbatim}
Example:
\begin{verbatim}
f.str("Hello")           #c# Returns "Hello"
f.str(42)                #c# Returns "42"
f.str("Value: ", 3.14)   #c# Returns "Value: 3.14"
\end{verbatim}

\subsubsection{String Formatting}
The \texttt{format} function formats a non-string datatype. \\
Morphology:
\begin{verbatim}
EBNF:
"f.format", "(", argument, [ ",", options ], ")"
\end{verbatim}
% TODO: Document specific format options

\subsubsection{String Case Functions}
\paragraph{Uppercase}
The \texttt{upperc} function converts a string to uppercase. \\
Morphology:
\begin{verbatim}
EBNF:
"f.upperc", "(", string, ")"
\end{verbatim}
Example:
\begin{verbatim}
f.upperc("hello")   #c# Returns "HELLO"
\end{verbatim}

\paragraph{Lowercase}
The \texttt{lowerc} function converts a string to lowercase. \\
Morphology:
\begin{verbatim}
EBNF:
"f.lowerc", "(", string, ")"
\end{verbatim}
Example:
\begin{verbatim}
f.lowerc("HELLO")   #c# Returns "hello"
\end{verbatim}


\subsection{List and Set Functions}

\subsubsection{Concatenation}
The \texttt{concat} function combines all arguments into a single list. \\
Morphology:
\begin{verbatim}
EBNF:
"f.concat", "(", { argument }, ")"
\end{verbatim}
Example:
\begin{verbatim}
f.concat([1, 2], [3, 4])        #c# Returns [1, 2, 3, 4]
f.concat("a", "b", "c")         #c# Returns ["a", "b", "c"]
\end{verbatim}

\subsubsection{Type Conversion}
\paragraph{List to Set}
The \texttt{set} function converts a list to a set. Only objects allowed in sets (parts, connections, pins, nets) are preserved; forbidden objects are silently dropped with a warning. \\
Morphology:
\begin{verbatim}
EBNF:
"f.set", "(", list, ")"
\end{verbatim}
Example:
\begin{verbatim}
f.set([Q1, Q2, R5])   #c# Returns {Q1, Q2, R5}
\end{verbatim}

\paragraph{Set to List}
The \texttt{list} function converts a set to a list. The order is determined by the order elements appear in \texttt{*e}. \\
Morphology:
\begin{verbatim}
EBNF:
"f.list", "(", set, ")"
\end{verbatim}
Example:
\begin{verbatim}
f.list({Q1, Q2})   #c# Returns [Q1, Q2] (order from *e)
\end{verbatim}


\subsection{Ndarray Functions}
Functions for operating on ndarrays (and lists). These functions extend the built-in math and list functions with array-aware behavior.

\subsubsection{Reduction Functions}
Reduction functions collapse an ndarray along an axis. When \texttt{axis} is omitted, the operation is applied to the flattened array (all elements). When \texttt{axis = 0}, the operation is applied row-wise (collapses rows). When \texttt{axis = 1}, the operation is applied column-wise (collapses columns).

\paragraph{Sum}
\begin{verbatim}
EBNF:
"f.sum", "(", argument, [ ",", "axis", "=", int ], ")"
\end{verbatim}
\begin{verbatim}
f.sum([1, 2, 3])              #c# Returns 6
f.sum([[1,2],[3,4]])          #c# Returns 10 (flattened)
f.sum([[1,2],[3,4]], axis=0)  #c# Returns [4, 6]  (row-wise)
f.sum([[1,2],[3,4]], axis=1)  #c# Returns [3, 7]  (column-wise)
\end{verbatim}

\paragraph{Product}
\begin{verbatim}
EBNF:
"f.prod", "(", argument, [ ",", "axis", "=", int ], ")"
\end{verbatim}
\begin{verbatim}
f.prod([2, 3, 4])              #c# Returns 24
f.prod([[2,3],[4,5]], axis=0)  #c# Returns [8, 15]
f.prod([[2,3],[4,5]], axis=1)  #c# Returns [6, 20]
\end{verbatim}

\paragraph{Mean}
\begin{verbatim}
EBNF:
"f.mean", "(", argument, [ ",", "axis", "=", int ], ")"
\end{verbatim}
\begin{verbatim}
f.mean([1, 2, 3, 4])            #c# Returns 2.5
f.mean([[1,2],[3,4]], axis=0)   #c# Returns [2.0, 3.0]
f.mean([[1,2],[3,4]], axis=1)   #c# Returns [1.5, 3.5]
\end{verbatim}

\paragraph{Standard Deviation}
\begin{verbatim}
EBNF:
"f.std", "(", argument, [ ",", "axis", "=", int ], ")"
\end{verbatim}
\begin{verbatim}
f.std([1, 2, 3, 4, 5])         #c# Returns ~1.4142
f.std([1, 1, 1, 1])            #c# Returns 0.0
\end{verbatim}

\paragraph{Variance}
\begin{verbatim}
EBNF:
"f.var", "(", argument, [ ",", "axis", "=", int ], ")"
\end{verbatim}
\begin{verbatim}
f.var([1, 2, 3, 4, 5])         #c# Returns 2.0
f.var([1, 1, 1, 1])            #c# Returns 0.0
\end{verbatim}

\paragraph{Median}
\begin{verbatim}
EBNF:
"f.median", "(", argument, ")"
\end{verbatim}
\begin{verbatim}
f.median([3, 1, 4, 1, 5])   #c# Returns 3
f.median([1, 2, 3, 4])      #c# Returns 2.5 (average of middle two)
\end{verbatim}

\paragraph{Percentile}
\begin{verbatim}
EBNF:
"f.percentile", "(", argument, ",", real, ")"
\end{verbatim}
The second argument is a percentage (0--100).
\begin{verbatim}
f.percentile([10, 20, 30, 40, 50], 50)  #c# Returns 30
f.percentile([10, 20, 30, 40, 50], 0)   #c# Returns 10
f.percentile([10, 20, 30, 40, 50], 100) #c# Returns 50
\end{verbatim}

\paragraph{Root Mean Square}
\begin{verbatim}
EBNF:
"f.rms", "(", argument, ")"
\end{verbatim}
Computes $\sqrt{\frac{1}{n}\sum_{i=1}^{n} x_i^2}$. Useful for AC signal analysis and power calculations.
\begin{verbatim}
f.rms([1, -1, 1, -1])        #c# Returns 1.0
f.rms([3, 4])                 #c# Returns 3.5355...
f.rms([0, 0, 0])             #c# Returns 0.0
\end{verbatim}

\paragraph{Peak}
\begin{verbatim}
EBNF:
"f.peak", "(", argument, [ ",", "mode", "=", string ], ")"
\end{verbatim}
Modes:
\begin{itemize}
\item \texttt{"p2p"} (default) -- Peak-to-peak: $\max(x) - \min(x)$
\item \texttt{"abs"} -- Maximum absolute value: $\max(|x|)$
\item \texttt{"pos"} -- Maximum positive value: $\max(x)$
\item \texttt{"neg"} -- Minimum negative value: $\min(x)$
\end{itemize}
\begin{verbatim}
f.peak([1, -3, 2, -5, 4])                  #c# Returns 9 (p2p: 4 - (-5))
f.peak([1, -3, 2, -5, 4], mode="f.abs")      #c# Returns 5
f.peak([1, -3, 2, -5, 4], mode="pos")      #c# Returns 4
f.peak([1, -3, 2, -5, 4], mode="neg")      #c# Returns -5
\end{verbatim}

\subsubsection{Search Functions}
\paragraph{Index of Minimum}
\begin{verbatim}
EBNF:
"f.argmin", "(", argument, ")"
\end{verbatim}
\begin{verbatim}
f.argmin([5, 3, 1, 4])   #c# Returns 2 (index of 1)
\end{verbatim}

\paragraph{Index of Maximum}
\begin{verbatim}
EBNF:
"f.argmax", "(", argument, ")"
\end{verbatim}
\begin{verbatim}
f.argmax([5, 3, 1, 4])   #c# Returns 0 (index of 5)
\end{verbatim}

\subsubsection{Shape Functions}
\paragraph{Flatten}
\begin{verbatim}
EBNF:
"f.flatten", "(", argument, ")"
\end{verbatim}
\begin{verbatim}
f.flatten([[1,2],[3,4],[5,6]])   #c# Returns [1, 2, 3, 4, 5, 6]
f.flatten([[1,2,3]])             #c# Returns [1, 2, 3]
\end{verbatim}

\paragraph{Sort}
\begin{verbatim}
EBNF:
"f.sort", "(", argument, [ ",", lambda ], ")"
\end{verbatim}
Returns a new sorted list (does not modify the original). An optional lambda comparator can be provided for custom ordering (see Higher-Order Functions section).
\begin{verbatim}
f.sort([5, 3, 1, 4])           #c# Returns [1, 3, 4, 5]
f.sort([3.1, 2.7, 1.4])       #c# Returns [1.4, 2.7, 3.1]
f.sort([5, 1, 3], L::($a [int], $b [int])->{ $b - $a })
#c# Returns [5, 3, 1]  (descending)
\end{verbatim}

\paragraph{Reverse}
\begin{verbatim}
EBNF:
"f.reverse", "(", argument, ")"
\end{verbatim}
Returns a new reversed list (does not modify the original).
\begin{verbatim}
f.reverse([1, 2, 3])           #c# Returns [3, 2, 1]
f.reverse([[1,2],[3,4]])       #c# Returns [[3,4],[1,2]]
\end{verbatim}

\paragraph{Unique}
\begin{verbatim}
EBNF:
"f.unique", "(", argument, ")"
\end{verbatim}
\begin{verbatim}
f.unique([1, 2, 2, 3, 3, 3])   #c# Returns [1, 2, 3]
\end{verbatim}

\subsubsection{Linear Algebra Functions}
\paragraph{Determinant}
\begin{verbatim}
EBNF:
"f.det", "(", argument, ")"
\end{verbatim}
\begin{verbatim}
f.det([[1,2],[3,4]])   #c# Returns -2.0
f.det([[2,0],[0,3]])   #c# Returns 6.0
\end{verbatim}

\paragraph{Matrix Inverse}
\begin{verbatim}
EBNF:
"f.inv", "(", argument, ")"
\end{verbatim}
Returns the inverse matrix as a list of lists. A Division by Zero Error is raised if the matrix is singular.
\begin{verbatim}
f.inv([[1,2],[3,4]])   #c# Returns [[-2.0, 1.0],[1.5, -0.5]]
\end{verbatim}

\paragraph{Eigenvalues}
\begin{verbatim}
EBNF:
"f.eig", "(", argument, ")"
\end{verbatim}
Returns a list of eigenvalues.
\begin{verbatim}
f.eig([[2,0],[0,3]])   #c# Returns [2, 3]
f.eig([[1,2],[2,1]])   #c# Returns [3, -1]
\end{verbatim}

\paragraph{Eigenvectors}
\begin{verbatim}
EBNF:
"f.eigvec", "(", argument, ")"
\end{verbatim}
Returns a list of eigenvectors (one per eigenvalue, as rows).
\begin{verbatim}
f.eigvec([[2,0],[0,3]])   #c# Returns [[1,0],[0,1]]
\end{verbatim}

\paragraph{Singular Value Decomposition}
\begin{verbatim}
EBNF:
"f.svd", "(", argument, ")"
\end{verbatim}
Returns a list of 3 lists: \texttt{[U, S, V]} where \texttt{U} and \texttt{V} are orthogonal matrices and \texttt{S} is the list of singular values.
\begin{verbatim}
f.svd([[1,0],[0,2]])
    #c# Returns [[[1,0],[0,1]], [1, 2], [[1,0],[0,1]]]
\end{verbatim}

\paragraph{LU Decomposition}
\begin{verbatim}
EBNF:
"f.lu", "(", argument, ")"
\end{verbatim}
Returns a list of 2 lists: \texttt{[L, U]} where \texttt{L} is lower triangular and \texttt{U} is upper triangular.
\begin{verbatim}
f.lu([[2,1],[4,3]])   #c# Returns L and U such that L .* U = original
\end{verbatim}

\paragraph{QR Decomposition}
\begin{verbatim}
EBNF:
"f.qr", "(", argument, ")"
\end{verbatim}
Returns a list of 2 lists: \texttt{[Q, R]} where \texttt{Q} is orthogonal and \texttt{R} is upper triangular.
\begin{verbatim}
f.qr([[1,1],[1,-1]])   #c# Returns Q and R such that Q .* R = original
\end{verbatim}

\paragraph{Cholesky Decomposition}
\begin{verbatim}
EBNF:
"f.chol", "(", argument, ")"
\end{verbatim}
Returns the lower triangular matrix \texttt{L} such that \texttt{L .* L.T = original}. A Value Error is raised if the matrix is not positive definite.
\begin{verbatim}
f.chol([[4,2],[2,3]])   #c# Returns [[2,0],[1,1.7321]]
\end{verbatim}

\paragraph{Trace}
\begin{verbatim}
EBNF:
"f.trace", "(", argument, ")"
\end{verbatim}
\begin{verbatim}
f.trace([[1,2],[3,4]])   #c# Returns 5
f.trace([[2,0],[0,3]])   #c# Returns 5
\end{verbatim}

\paragraph{Matrix Rank}
\begin{verbatim}
EBNF:
"f.rank", "(", argument, ")"
\end{verbatim}
\begin{verbatim}
f.rank([[1,2],[3,4]])   #c# Returns 2
f.rank([[1,2],[2,4]])   #c# Returns 1 (rows are linearly dependent)
\end{verbatim}

\paragraph{Condition Number}
\begin{verbatim}
EBNF:
"f.cond", "(", argument, ")"
\end{verbatim}
Returns the ratio of the largest to smallest singular value. A high condition number indicates numerical instability.
\begin{verbatim}
f.cond([[1,0],[0,1]])   #c# Returns 1.0 (well-conditioned)
f.cond([[1,0],[0,1e-10]]) #c# Returns 1e10 (ill-conditioned)
\end{verbatim}

\paragraph{Pseudo-Inverse}
\begin{verbatim}
EBNF:
"f.pinv", "(", argument, ")"
\end{verbatim}
Returns the Moore--Penrose pseudo-inverse. Works for non-square or singular matrices.
\begin{verbatim}
f.pinv([[1,2],[3,4]])   #c# Returns [[-2, 1],[1.5, -0.5]]
\end{verbatim}

\subsubsection{Statistical Functions}
\paragraph{Histogram}
\begin{verbatim}
EBNF:
"f.histogram", "(", argument, ",", int, ")"
\end{verbatim}
Returns a list of 2 lists: \texttt{[edges, counts]} where \texttt{edges} has \texttt{n + 1} values and \texttt{counts} has \texttt{n} values.
\begin{verbatim}
f.histogram([1, 2, 3, 4, 5], 3)
#c# Returns [[1.0, 2.333, 3.667, 5.0], [2, 1, 2]]
\end{verbatim}


\subsection{Higher-Order Functions}
Higher-order functions take lambda expressions as arguments. They operate on lists and ndarrays. See Section~\ref{sec:lambda} for lambda syntax.

\subsubsection{Map}
The \texttt{map} function applies a lambda to each element of a list and returns a new list with the results. The input list is not modified.

\paragraph{Morphology}
\begin{verbatim}
EBNF:
"f.map", "(", list, ",", lambda, ")"
\end{verbatim}

\paragraph{Type Overloading}
\begin{itemize}
\item \texttt{map(list[T], L::(\$x [T])->\{...\})} $\rightarrow$ \texttt{list[U]}
  (T and U can be different types)
\item Works on any list: flat, nested, mixed types
\item The lambda parameter type must match the element type
\end{itemize}

\paragraph{Examples}
\begin{verbatim}
f.map([1, 2, 3], L::($x [int])->{ $x * 2 })
#c# Returns [2, 4, 6]

f.map([1, 2, 3], L::($x [int])->{ $x > 2 })
#c# Returns [False, False, True]  (type change: int -> bool)

f.map(["a", "bb", "ccc"], L::($x [string])->{ f.len($x) })
#c# Returns [1, 2, 3]  (type change: string -> int)

f.map([1, 2, 3], L::($x [int])->{
    $y [int]= $x * 10
    $y + 1
})
#c# Returns [11, 21, 31]  (multi-line lambda)

f.map([], L::($x [int])->{ $x * 2 })
#c# Returns []  (empty list)
\end{verbatim}

\paragraph{Error Conditions}
\begin{itemize}
\item Type mismatch: lambda parameter type does not match element type raises a Runtime Error
\end{itemize}

\subsubsection{Filter}
The \texttt{filter} function returns a new list containing only elements for which the lambda returns \texttt{True}. The input list is not modified.

\paragraph{Morphology}
\begin{verbatim}
EBNF:
"f.filter", "(", list, ",", lambda, ")"
\end{verbatim}

\paragraph{Type Overloading}
\begin{itemize}
\item \texttt{filter(list[T], L::(\$x [T])->\{bool\})} $\rightarrow$ \texttt{list[T]}
  (return type is the same as input)
\item The lambda must return \texttt{bool}
\end{itemize}

\paragraph{Examples}
\begin{verbatim}
f.filter([1, 2, 3, 4, 5], L::($x [int])->{ $x % 2 == 0 })
#c# Returns [2, 4]

f.filter([1, 2, 3, 4, 5], L::($x [int])->{ $x > 10 })
#c# Returns []  (no matches)

f.filter(["a", "bb", "ccc", "dd"], L::($x [string])->{ f.len($x) > 2 })
#c# Returns ["ccc"]

f.filter([-1, -2, 3, -4, 5], L::($x [int])->{ $x > 0 })
#c# Returns [3, 5]

f.filter([1, 2, 3], L::($x [int])->{ $x > 5 })
#c# Returns []  (empty result)

f.filter([], L::($x [int])->{ $x > 0 })
#c# Returns []  (empty input)
\end{verbatim}

\paragraph{Error Conditions}
\begin{itemize}
\item Type mismatch: lambda parameter type does not match element type raises a Runtime Error
\item Lambda does not return \texttt{bool} raises a Type Error
\end{itemize}

\subsubsection{Reduce}
The \texttt{reduce} function folds a list to a single value by applying a binary lambda cumulatively. The lambda receives an accumulator and the current element. The \texttt{init} value is always required.

\paragraph{Morphology}
\begin{verbatim}
EBNF:
"f.reduce", "(", list, ",", lambda, ",", expression, ")"
\end{verbatim}

\paragraph{Type Overloading}
\begin{itemize}
\item \texttt{reduce(list[T], L::(\$acc [U], \$x [T])->\{U\}, init[U])} $\rightarrow$ \texttt{U}
  (accumulator type U can differ from element type T)
\item The \texttt{init} value type must match the accumulator type
\end{itemize}

\paragraph{Examples}
\begin{verbatim}
#c# Sum integers
f.reduce([1, 2, 3, 4], L::($acc [int], $x [int])->{ $acc + $x }, 0)
#c# Returns 10

#c# Product integers
f.reduce([2, 3, 4], L::($acc [int], $x [int])->{ $acc * $x }, 1)
#c# Returns 24

#c# Find maximum
f.reduce([3, 1, 4, 1, 5], L::($acc [int], $x [int])->{ 
    ($acc > $x)?->($acc)!($x) 
}, 0)
#c# Returns 5

#c# Reduce with string concatenation
f.reduce(["a", "b", "c"],
    L::($acc [string], $x [string])->{ $acc + $x }, "")
#c# Returns "abc"

#c# Reduce on empty list returns init
f.reduce([], L::($acc [int], $x [int])->{ $acc + $x }, 42)
#c# Returns 42

#c# Flatten with reduce (using concat)
f.reduce([[1,2],[3,4]],
    L::($acc [list], $x [list])->{ f.concat($acc, $x) }, [])
#c# Returns [1, 2, 3, 4]
\end{verbatim}

\paragraph{Error Conditions}
\begin{itemize}
\item Type mismatch: lambda parameter types do not match accumulator/element types raises a Runtime Error
\item Lambda does not return a type coercible to the accumulator type raises a Type Error
\end{itemize}

\subsubsection{Where}
The \texttt{where} function returns a list of \textbf{indices} (ints) for which the lambda returns \texttt{True}. Indices are 0-based.

\paragraph{Morphology}
\begin{verbatim}
EBNF:
"f.where", "(", list, ",", lambda, ")"
\end{verbatim}

\paragraph{Type Overloading}
\begin{itemize}
\item \texttt{where(list[T], L::(\$x [T])->\{bool\})} $\rightarrow$ \texttt{list[int]}
\item The lambda must return \texttt{bool}
\end{itemize}

\paragraph{Examples}
\begin{verbatim}
f.where([-1, 3, -2, 5], L::($x [int])->{ $x > 0 })
#c# Returns [1, 3]

f.where([10, 20, 30], L::($x [int])->{ $x >= 20 })
#c# Returns [1, 2]

f.where([1, 2, 3], L::($x [int])->{ $x > 10 })
#c# Returns []  (no matches)

f.where(["a", "bb", "ccc"], L::($x [string])->{ f.len($x) >= 2 })
#c# Returns [1, 2]

f.where([], L::($x [int])->{ $x > 0 })
#c# Returns []  (empty input)
\end{verbatim}

\paragraph{Error Conditions}
\begin{itemize}
\item Type mismatch: lambda parameter type does not match element type raises a Runtime Error
\item Lambda does not return \texttt{bool} raises a Type Error
\end{itemize}

\subsubsection{Count}
The \texttt{count} function returns the number of elements for which the lambda returns \texttt{True}.

\paragraph{Morphology}
\begin{verbatim}
EBNF:
"f.count", "(", list, ",", lambda, ")"
\end{verbatim}

\paragraph{Type Overloading}
\begin{itemize}
\item \texttt{count(list[T], L::(\$x [T])->\{bool\})} $\rightarrow$ \texttt{int}
\item The lambda must return \texttt{bool}
\end{itemize}

\paragraph{Examples}
\begin{verbatim}
f.count([1, 2, 3, 4, 5], L::($x [int])->{ $x > 3 })
#c# Returns 2

f.count([True, False, True, True], L::($x [bool])->{ $x })
#c# Returns 3

f.count([1, 2, 3], L::($x [int])->{ $x > 10 })
#c# Returns 0

f.count([], L::($x [int])->{ $x > 0 })
#c# Returns 0
\end{verbatim}

\paragraph{Error Conditions}
\begin{itemize}
\item Type mismatch: lambda parameter type does not match element type raises a Runtime Error
\item Lambda does not return \texttt{bool} raises a Type Error
\end{itemize}

\subsubsection{Find}
The \texttt{find} function returns the first element for which the lambda returns \texttt{True}. If no element matches, a Value Error is raised.

\paragraph{Morphology}
\begin{verbatim}
EBNF:
"f.find", "(", list, ",", lambda, ")"
\end{verbatim}

\paragraph{Type Overloading}
\begin{itemize}
\item \texttt{find(list[T], L::(\$x [T])->\{bool\})} $\rightarrow$ \texttt{T}
  (return type is the same as element type)
\item The lambda must return \texttt{bool}
\end{itemize}

\paragraph{Examples}
\begin{verbatim}
f.find([1, 2, 3, 4, 5], L::($x [int])->{ $x > 2 })
#c# Returns 3

f.find([1, 2, 3, 4, 5], L::($x [int])->{ $x == 4 })
#c# Returns 4

f.find(["a", "bb", "ccc"], L::($x [string])->{ f.len($x) > 1 })
#c# Returns "bb"

f.find([1, 2, 3, 4, 5], L::($x [int])->{ $x > 10 })
#c# Value Error: no matching element

f.find([], L::($x [int])->{ $x > 0 })
#c# Value Error: no matching element
\end{verbatim}

\paragraph{Error Conditions}
\begin{itemize}
\item Value Error: no element matches the lambda condition
\item Type mismatch: lambda parameter type does not match element type raises a Runtime Error
\item Lambda does not return \texttt{bool} raises a Type Error
\end{itemize}

\subsubsection{Sort with Comparator}
The \texttt{sort} function returns a new sorted list. Without a comparator, default sort order is ascending. An optional lambda comparator can be provided for custom ordering.

\paragraph{Morphology}
\begin{verbatim}
EBNF:
"f.sort", "(", list, [ ",", lambda ], ")"
\end{verbatim}

\paragraph{Comparator Lambda}
The comparator lambda receives two elements (\texttt{\$a} and \texttt{\$b}) and must return a numeric value:
\begin{itemize}
\item \textbf{Negative} value (e.g., \texttt{\$a - \$b}): \texttt{\$a} comes before \texttt{\$b}
\item \textbf{Zero}: order unchanged (stable sort)
\item \textbf{Positive} value (e.g., \texttt{\$b - \$a}): \texttt{\$b} comes before \texttt{\$a}
\end{itemize}

This is equivalent to the \texttt{cmp} convention used in many languages. The simplest ascending comparator is \texttt{\$a - \$b}, and descending is \texttt{\$b - \$a}.

\paragraph{Type Overloading}
\begin{itemize}
\item \texttt{sort(list[T])} $\rightarrow$ \texttt{list[T]} (default ascending)
\item \texttt{sort(list[T], L::(\$a [T], \$b [T])->\{int\})} $\rightarrow$ \texttt{list[T]}
  (custom comparator, must return \texttt{int} or \texttt{real})
\end{itemize}

\paragraph{Examples}
\begin{verbatim}
#c# Default sort (ascending)
f.sort([5, 1, 3])
#c# Returns [1, 3, 5]

#c# Ascending with comparator
f.sort([5, 1, 3], L::($a [int], $b [int])->{ $a - $b })
#c# Returns [1, 3, 5]

#c# Descending with comparator
f.sort([5, 1, 3], L::($a [int], $b [int])->{ $b - $a })
#c# Returns [5, 3, 1]

#c# Sort strings by length
f.sort(["bb", "a", "ccc"],
    L::($a [string], $b [string])->{ f.len($a) - f.len($b) })
#c# Returns ["a", "bb", "ccc"]

#c# Sort reals ascending
f.sort([3.14, 1.41, 2.72], L::($a [real], $b [real])->{ $a - $b })
#c# Returns [1.41, 2.72, 3.14]

#c# Sort with nested sort key (sort by absolute value)
f.sort([-5, 3, -1, 4], L::($a [int], $b [int])->{ f.abs($a) - f.abs($b) })
#c# Returns [-1, 3, 4, -5]

#c# Sort preserves original (returns new list)
$list [list]= [5, 3, 1]
$sorted [list]= f.sort($list)
#c# $list is still [5, 3, 1], $sorted is [1, 3, 5]
\end{verbatim}

\paragraph{Error Conditions}
\begin{itemize}
\item Type mismatch: comparator parameter types do not match element types raises a Runtime Error
\item Comparator does not return a numeric type raises a Type Error
\end{itemize}

\subsection{Signal Processing Functions}
Functions for frequency-domain analysis of signals.

\subsubsection{Fast Fourier Transform}
The \texttt{fft} function computes the Discrete Fourier Transform of a real-valued signal.
\begin{verbatim}
EBNF:
"f.fft", "(", list, ")"
\end{verbatim}

\paragraph{Type Overloading}
\begin{itemize}
\item \texttt{fft(list[int])} $\rightarrow$ \texttt{list[complex]}
\item \texttt{fft(list[real])} $\rightarrow$ \texttt{list[complex]}
\item Output length equals input length
\item The input signal is zero-padded to the next power of 2 if needed
\end{itemize}

\paragraph{Examples}
\begin{verbatim}
#c# FFT of a simple signal
f.fft([1, 0, 1, 0])   #c# Returns [2+0i, 0+0i, 2+0i, 0+0i]

#c# FFT of a single impulse
f.fft([1, 0, 0, 0])   #c# Returns [1+0i, 1+0i, 1+0i, 1+0i]

#c# FFT of a DC signal (constant value)
f.fft([5, 5, 5, 5])   #c# Returns [20+0i, 0+0i, 0+0i, 0+0i]

#c# Extract magnitude spectrum
$freq [list]= f.fft([1, 0, 1, 0])
$mags [list]= f.map($freq, L::($z [complex])->{ f.abs($z) })
#c# $mags = [2.0, 0.0, 2.0, 0.0]
\end{verbatim}

\paragraph{Error Conditions}
\begin{itemize}
\item Empty list raises a Value Error
\end{itemize}

\subsubsection{Inverse Fast Fourier Transform}
The \texttt{ifft} function computes the Inverse Discrete Fourier Transform, reconstructing a real signal from complex frequency components.
\begin{verbatim}
EBNF:
"f.ifft", "(", list, ")"
\end{verbatim}

\paragraph{Type Overloading}
\begin{itemize}
\item \texttt{ifft(list[complex])} $\rightarrow$ \texttt{list[real]}
\item Output length equals input length
\item Imaginary parts are discarded (reconstructed signal is real)
\end{itemize}

\paragraph{Examples}
\begin{verbatim}
#c# Round-trip: fft then ifft
$signal [list]= [1, 0, 1, 0]
$freq [list]= f.fft($signal)
$reconstructed [list]= f.ifft($freq)
#c# $reconstructed is approximately [1, 0, 1, 0]

#c# ifft of DC component
f.ifft([4+0i, 0+0i, 0+0i, 0+0i])   #c# Returns [1, 1, 1, 1]

#c# ifft of impulse in frequency domain
f.ifft([1+0i, 1+0i, 1+0i, 1+0i])   #c# Returns [1, 0, 0, 0]
\end{verbatim}

\paragraph{Error Conditions}
\begin{itemize}
\item Empty list raises a Value Error
\end{itemize}

\subsection{Shared Functions}

\subsubsection{Length/Size Function}
The \texttt{len} function returns the length of a list or string, or the size of a set. \\
Morphology:
\begin{verbatim}
EBNF:
"f.len", "(", ( list | string | set ), ")"
\end{verbatim}
Examples:
\begin{verbatim}
f.len([1, 2, 3])     #c# Returns 3
f.len("hello")       #c# Returns 5
f.len({Q1, Q2, R5})  #c# Returns 3
\end{verbatim}


\subsection{Special Functions}

\subsubsection{Print Function}
The \texttt{f.print()} function outputs its argument to the console. Unlike \texttt{echo}, it can evaluate expressions and function calls directly. \\
Morphology:
\begin{verbatim}
EBNF:
"f.print", "(", argument, ")"
\end{verbatim}
Example:
\begin{verbatim}
f.print("Hello, World!")
f.print(3 + 2 * 5)   #c# Outputs: 13
\end{verbatim}

\subsubsection{Type Function}
The \texttt{type} function returns the type of an object as a string. \\
Morphology:
\begin{verbatim}
EBNF:
"f.type", "(", argument, ")"
\end{verbatim}
Example:
\begin{verbatim}
f.type(42)       #c# Returns "int"
f.type(3.14)     #c# Returns "real"
f.type("hello")  #c# Returns "string"
f.type(Q1)       #c# Returns "part"
\end{verbatim}

\subsubsection{Reference Function}
The \texttt{ref} function converts a string containing a part reference into the referenced part. \\
Morphology:
\begin{verbatim}
EBNF:
"f.ref", "(", string, ")"
\end{verbatim}
Example:
\begin{verbatim}
f.ref("R1")   #c# Returns the part referenced as R1
\end{verbatim}

\subsubsection{No-Operation Function}
The \texttt{noop} function is the identity function. It returns its argument unchanged. \\
Morphology:
\begin{verbatim}
EBNF:
"f.noop", "(", argument, ")"
\end{verbatim}
Example:
\begin{verbatim}
f.noop(42)       #c# Returns 42
f.noop("hello")  #c# Returns "hello"
f.noop(Q1)       #c# Returns Q1
\end{verbatim}
Use \texttt{noop} as a placeholder in queries where no transformation is needed.

\subsubsection{Soft Integer Conversion}
The \texttt{softint} function converts its argument to an integer if it is either already an integer or is a real number numerically equal to an integer. Otherwise it returns the argument unchanged. \\
Morphology:
\begin{verbatim}
EBNF:
"f.softint", "(", argument, ")"
\end{verbatim}
Examples:
\begin{verbatim}
f.softint(42)       #c# Returns 42 (already int)
f.softint(5.0)      #c# Returns 5 (real equal to int)
f.softint(3.7)      #c# Returns 3.7 (not equal to int, unchanged)
f.softint("hello")  #c# Returns "hello" (not numeric, unchanged)
\end{verbatim}

\subsubsection{Soft Real Conversion}
The \texttt{softreal} function converts its argument to a real if it is either already a real, an integer, or a complex number with zero imaginary part. Otherwise it returns the argument unchanged. \\
Morphology:
\begin{verbatim}
EBNF:
"f.softreal", "(", argument, ")"
\end{verbatim}
Examples:
\begin{verbatim}
f.softreal(42)        #c# Returns 42.0 (int promoted to real)
f.softreal(3.14)      #c# Returns 3.14 (already real)
f.softreal(5 + 0i)    #c# Returns 5.0 (complex with zero imag part)
f.softreal(3 + 4i)
    #c# Returns 3 + 4i (complex with non-zero imag, unchanged)
f.softreal("hello")   #c# Returns "hello" (not numeric, unchanged)
\end{verbatim}

\section{User-defined Functions}
User-defined functions allow creation of reusable code blocks with local scoping.

\subsection{Function Definition}
Functions are defined using the \texttt{function} keyword with the \texttt{f.u.} prefix:
\begin{verbatim}
EBNF:
function_definition = "function", "f.u.", id, 
  "(", [ { id, "," }, id], ")", "::", "{", 
  f_statement, {f_statement}, "}"

f_statement = statement | "return", [data]

Example:
function f.u.add( $a, $b )::{
    $retval [int]= $a + $b
}
\end{verbatim}

Function parameters are local variables scoped to the function body. To return a value, set the \texttt{\$retval} variable.

\subsection{Function Calls}
User-defined functions are called using the \texttt{f.u.} prefix:
\begin{verbatim}
EBNF:
function_call_ud = "f.u.", id, "(", 
                   [{ data, "," }, data], ")"

Example:
$result [int]= f.u.add( 3, 5 )   #c# $result = 8
\end{verbatim}

\subsection{Return Values}
Functions return values by setting the \texttt{\$retval} variable. If \texttt{\$retval} is not set, the function returns nothing:
\begin{verbatim}
function f.u.square( $x )::{
    $retval [int]= $x ^ 2
}

function f.u.greet( $name )::{
    f.print( f.str( "Hello, ", $name, "!" ) )
    #c# No return value
}

$sq [int]= f.u.square( 4 )     #c# $sq = 16
f.u.greet( "World" )          #c# Prints: Hello, World!
\end{verbatim}

\subsection{Local Scoping}
\label{sec:function-scope}
Functions have local scoping. Variables declared inside a function are only accessible within that function. Parameters are also local variables.

\paragraph{Scope Rules}
\begin{itemize}
\item Parameters (\$a, \$b, etc.) are local to the function body
\item Variables declared inside the function with \texttt{\$var [type]= val} are local
\item Local variables shadow global variables of the same name
\item The \texttt{\$retval} variable is always local to the function
\item Functions cannot access local variables from other function invocations
\end{itemize}

\paragraph{Return Values}
Functions can return values using two equivalent mechanisms:
\begin{itemize}
\item Assign to \texttt{\$retval} and let execution fall through the closing \texttt{\}}
\item Use the \texttt{return} keyword, which is syntactic sugar that assigns its argument to \texttt{\$retval} and immediately exits the function
\end{itemize}
Both \texttt{return \$value} and \texttt{\$retval [type]= \$value} followed by natural function exit are valid and equivalent.

\paragraph{Examples}
\begin{verbatim}
$x [int]= 10      #c# Global variable

function f.u.doublen()::{
    $x [int]= $x * 2  #c# Local $x shadows global
    f.print($x)            #c# Prints: 20
}

f.u.doublen()
f.print($x)          #c# Prints: 10 (global unchanged)

function f.u.addLocal( $a, $b )::{
    $sum [int]= $a + $b   #c# $sum is local
    $retval [int]= $sum
}

f.u.addLocal( 3, 4 )   #c# Returns 7
#c# $sum is not accessible here
\end{verbatim}

\paragraph{Recursion}
Local scoping enables recursion since each call has its own copy of parameters:
\begin{verbatim}
function f.u.factorial( $n )::{
    if $n <= 1 ::{
        $retval [int]= 1
    } else ::{
        $retval [int]= $n * 
            f.u.factorial( $n - 1 )
    }
}

$result [int]= f.u.factorial( 5 )  #c# = 120
\end{verbatim}




\part{Formal Semantics}
